\section{Related work}

\textbf{Home-Cage modelling.}
HCM is an active and rapidly growing field, with recent surveys describing available systems, sensing technologies, experimental applications, and do-it-yourself approaches \cite{kiryk2026,huzard2026tech,duarte2026diy}. In parallel, the HCM community has increasingly addressed the challenges created by heterogeneous data formats, incomplete metadata, and limited interoperability. Bains et al.\ \cite{bains2025} proposed a FAIR-compliant repository architecture for HCM data and metadata, emphasising standardised data structures, comprehensive metadata, and ontology-supported harmonisation as prerequisites for cross-system data reuse. This direction was subsequently developed further in work specifically addressing HCM data sharing and metadata \cite{forrest2026}, and connected to the broader ethical argument that preserving the scientific value and reusability of animal-derived data is itself part of responsible animal research \cite{petit2026wellfair}.

Importantly, semantic modelling of HCM also predates the present ontology. Within COST Action TEATIME, an Olog-based representation was developed to formalise what constitutes a home-cage monitoring system \cite{teatimeHCMDefinition,kiryk2026}. The work, led by Leonardo Restivo with contributions from Davor Virag and other TEATIME members, represented relationships among the animal, enclosure, hardware, software, sensors and actuators, behaviour and physiology, environmental provisions, duration of monitoring, and limited human interaction. The model was intended both to distinguish HCM systems from other experimental configurations and to identify key parameters that should be reported to describe an HCM experiment. This conceptual structure was also incorporated into the FAIR HCM repository work of Bains et al.\ \cite{bains2025}, providing an important precursor to the present work.

HCMO draws on this HCM-specific conceptual foundation while addressing a different level of formalisation. The TEATIME Olog provides a graphical, category-inspired conceptual model and explicitly adopts a simplified Olog representation rather than a deployable semantic-web ontology \cite{teatimeHCMDefinition}. Similarly, previous HCM data-sharing work establishes the need for standardised metadata, controlled terminology, and interoperable data structures, but does not provide a shared ontology capable of representing HCM experimental entities and their relationships in a machine-actionable form. HCMO addresses this formalisation gap through an explicit semantic model of the HCM acquisition chain—linking animal subjects, housing and environmental context, devices and sensors, observations, derived measurements, behaviours, and results—with formally defined entities and relations intended for annotation, validation, integration, and computational reuse across heterogeneous HCM platforms.


\textbf{Standards reused, not duplicated.} HCMO builds on W3C and community vocabularies. SOSA/SSN provides the sensor/observation/platform backbone \cite{sosa}; OWL-Time models temporal entities and intervals \cite{owltime}; QUDT provides quantity values and units \cite{qudt}; PROV captures provenance \cite{provo}; and BFO supplies an upper-ontology grounding for processes such as behaviour \cite{bfo}. SemTS provides a vocabulary for time-series segments and their dimensions \cite{semts}. HCMO reuses reviewed portions of these resources, contributing the HCM-specific concepts and relations that bind them.

\textbf{Adjacent biomedical and laboratory-animal ontologies.} Several resources cover neighbouring concerns. The Ontology for Biomedical Investigations (OBI) models investigations, protocols, instruments, materials, and data \cite{obi}; the Ontology of Laboratory Animal Medicine (OLAM) provides laboratory-animal terminology; and the Mouse Experimental Design Ontology (MEDO) standardises experimental-design descriptions, overlapping HCMO's housing/experimental-group notions. Reporting guidelines such as ARRIVE \cite{arrive} prescribe \emph{what} to report about in-vivo experiments but are not machine-readable ontologies. None of these model continuous in-cage acquisition; HCMO is complementary to them and can bridge to OBI and MEDO where appropriate.

The ISA framework provides a widely used investigation--study--assay exchange model \cite{isa}, while RO-Crate packages research data and contextual metadata as JSON-LD \cite{rocrate}. HCMO uses these as exchange targets in executable fixtures. This differs from asserting OWL mappings or complete formal conformance: the native ISA projections intentionally expose losses for HCM-specific housing, repeated-observation, and statistical-result structures.

\textbf{Ontology engineering.} Methodologically, HCMO follows an incremental, competency-question-driven process in the spirit of LOT \cite{lot2022} and SAMOD \cite{samod2017}, with requirements expressed as competency questions \cite{noy2001}. The current version was conceptualised graphically in diagrams.net and exported to OWL/Turtle with Chowlk \cite{chowlk2022}; LinkML was explored as an alternative serialisation \cite{linkml2021}. Documentation is produced with WIDOCO \cite{widoco2017}, and quality is assessed with OOPS! \cite{oops2012} and, for FAIRness, FOOPS! \cite{foops}, complemented by reasoning and SHACL validation \cite{gangemi2006}. Positioned this way, HCMO is the first ontology dedicated to HCM, designed to interoperate with --- rather than replace --- the surrounding standards and biomedical ontologies.
